\chapter{Basic Definitions}

\section{Clap}

A clap --- I prefer to call it a beat, but I use clap to avoid confusion of terms ---
is the pulse, and is called Akshara in Sanskrit. A cycle of aksharas is the Avartana
(bar). I like to think of this as being exactly the same as in Western Music. In
short:
\begin{align*}
  \text{Beat} &\sim \text{Akshara}\\
  \text{Bar}  &\sim \text{Avartana}
\end{align*}
However, I will refrain from using the term Beat. To avoid confusion, I will instead
use clap. A cycle of a particular number of claps is called the Tala.\footnote{Though
the definitions of Avartana and Tala seem to be the same, they are not. An avartana
is a measurement unit, whereas Tala is an identity, a name given to a particular
clapping pattern.} The definition of Tala extends beyond the number of claps --- it
also denotes the clapping pattern. The following sections strive to explain the
different Talas in carnatic music.

\subsection{Tala types and clapping}

There are different ways to categorize talas. However, the most fundamental talas are
below. In my opinion, this classification is the most fundamental one. There are
other classifications too. Please look at the wikipedia article for tala. These talas
are fundamental in the sense that learning these will enable a person to learn all
the other talas easily --- of course, you can express any number as a sum of these
numbers.

\begin{center}
\begin{adjustbox}{max width=\textwidth}
\begin{tabular}{c l c}
\toprule
S.No. & Tala name & Number of claps\\
\midrule
1 & Adi             & 8\\
2 & Rupakam         & 3 or 6\\
3 & Kanda Chapu     & 2.5\\
4 & Misra Chapu     & 3.5\\
5 & Sankeerna Chapu & 4.5\\
\bottomrule
\end{tabular}
\end{adjustbox}
\end{center}

To count or clap-out these talas we will need to understand certain terms.
\begin{enumerate}
  \item Double-clap or Drutam: Clap once, turn the right hand about, clap again.
  \item Single-clap or Anudrutam: Clap once.
  \item One-and-a-half clap: Clap twice at twice the speed of a single-clap, and give
    a pause for half a clap's duration. Thus, quite obviously, a one-and-a-half clap
    is one and a half claps long.
  \item Clap-count or Lagu: Clap once, on the second clap only the pinky touches the
    other hand, on the third clap only the ring-finger touches the other hand, on the
    fourth clap, only the middle-finger touches the other hand, and so on.
\end{enumerate}

The Clap-count can vary in length. They are usually denoted using sanskrit names for
the numbers. Chatusra is 4, Tisra is 3, Kanda is 5, Misra is 7, Sankeerna is 9. The
above tala table is presented below with the counting method included.

\begin{center}
\begin{adjustbox}{max width=\textwidth}
\begin{tabular}{c l c p{6.2cm}}
\toprule
S.No. & Tala name & Number of claps & Counting method\\
\midrule
1 & Adi             & 8   & 4-clap-count, double-clap, double-clap\\
2 & Rupakam         & 6   & double-clap, 4-clap-count\\
3 & Kanda Chapu     & 2.5 & single-clap, one-and-a-half clap\\
4 & Misra Chapu     & 3.5 & one-and-a-half clap with inverted palm, single-clap, single-clap\\
5 & Sankeerna Chapu & 4.5 & clap-count to 3 and then a one-and-a-half clap on the fourth count\\
\bottomrule
\end{tabular}
\end{adjustbox}
\end{center}

\section{Atom and Walk}

A Matra or atom is the most indivisible unit of time measurement in carnatic music.
Every clap/beat consists of a few or several atoms. The number of atoms per clap is
called walk or Nadai (Tamil for walk).

Walk decides how fast a piece is. For example, 'tha dhi ki ta thom' is always five
atoms long. If you are playing a 'tha dhi ki ta thom' at the speed of four atoms per
clap it would last for 1.25 claps. But, if you play it at five atoms per clap it
would last for 1 clap. The above statement makes one crucial assumption --- each
syllable in 'tha dhi ki ta thom' lasts for exactly the same duration. The same
assumption holds for the entire vocabulary of terms.

Note: Walk is never a fraction. It is always a whole number. This is because it is
tough to count if there are, for example, 4.5 atoms per clap. There are very few
exceptions to this rule --- only when it is easy to count.

A few nadais are named conveniently:

\begin{center}
\begin{adjustbox}{max width=\textwidth}
\begin{tabular}{l c}
\toprule
Nadai & atoms/clap\\
\midrule
Tisram     & 3\\
Chatusram  & 4\\
Kandam     & 5\\
Misram     & 7\\
Sankeernam & 9\\
\bottomrule
\end{tabular}
\end{adjustbox}
\end{center}
